\SourceLeaf{235}%
and consequently
\[
\begin{aligned}
\frac{\mathrm{d}X}{\mathrm{d}x}&=Z^3(-(1+uu)T+tuU)\\
\frac{\mathrm{d}X}{\mathrm{d}y}&=Z^3(-(1+uu)U+tuV)\\
\frac{\mathrm{d}Y}{\mathrm{d}x}&=Z^3(tuT-(1+tt)U)\\
\frac{\mathrm{d}Y}{\mathrm{d}y}&=Z^3(tuU-(1+tt)V)
\end{aligned}
\]
Substituting these values into the preceding expression gives
\[
k=Z^6(TV-UU)(1+tt+uu)=Z^4(TV-UU)=\frac{TV-UU}{(1+tt+uu)^2}
\]
\Article{8}
By a suitable choice of origin and coordinate axes, we can readily arrange that, at a given point \(A\), the values of \(t,u,U\) vanish. In fact, the first two conditions are already met if we take the tangent plane at this point as the plane of the coordinates \(x,y\). If we also place the origin of these coordinates at \(A\) itself, the expression for the coordinate \(z\) evidently takes the form
\[
z=\tfrac{1}{2}T^0xx+U^0xy+\tfrac{1}{2}V^0yy+\Omega
\]
where \(\Omega\) is of order greater than two. Then, on rotating the axes \(x,y\) through an angle \(M\) such that
\[
\mathrm{tang}\,2M=\frac{2U^0}{T^0-V^0}
\]
we readily see that the resulting equation has the form
\[
z=\tfrac{1}{2}Txx+\tfrac{1}{2}Vyy+\Omega
\]
The third condition is thus satisfied as well. Once this has been done, the following is clear.

I. If the curved surface is cut by a plane normal to it and passing through the coordinate axis \(x\), the result is a plane curve whose radius of curvature at \(A\) is \(=\frac{1}{T}\). The positive or negative sign indicates concavity or convexity toward the side in the direction of positive \(z\) coordinates.

II. Similarly, \(\frac{1}{V}\) is the radius of curvature at \(A\) of the plane curve formed by cutting the curved surface with the plane through the axes \(y,z\).
